\chapter{Implementación, despliegue y pruebas}
\label{chap:chapter3}

El presente capítulo materializa los esfuerzos de análisis y diseño expuestos en los capítulos anteriores,
trasladando las especificaciones técnicas a una implementación concreta y funcional. Aquí se detallan las
\ac{TI} ejecutadas para construir el servidor \textit{backend}, el diagrama de despliegue que ilustra la 
distribución física de los componentes del sistema en el entorno de ejecución , así como el proceso de 
pruebas al que fue sometido el sistema para garantizar su
calidad, rendimiento, seguridad y conformidad con los requisitos establecidos. La verificación y validación
sistemáticas mediante diversas estrategias y tipos de pruebas son fundamentales para asegurar que la 
solución desarrollada satisface las necesidades del cliente final y se encuentre lista para su
despliegue en un entorno productivo.


\section{Fase de implementación}

Dentro del ciclo de vida de un sistema informático se encuentra la fase de implementación. Esta
es la fase más costosa y que consume más tiempo. Se dice costosa ya que muchas personas,
herramientas y recursos están involucrados. La metodología \ac{XP} propone que las historias de
usuario deben ser implementadas en dependencia de la iteración a la que pertenezcan y las
mismas se descomponen en tareas de ingeniería \citep{maida2015}.

\subsection{Tareas de Ingeniería}

Las \ac{TI} constituyen un artefacto fundamental dentro de la metodología \ac{XP}. 
Representan la descomposición de las \ac{HU} en actividades técnicas concretas y accionables por parte de los
desarrolladores. Mientras que una \ac{HU} describe una funcionalidad desde la perspectiva del cliente, una \ac{TI}
detalla el “cómo” se implementará a nivel técnico \citep{pressman2014se}. En esta sección se enumeran y describen 
estas tareas, que guiaron el trabajo de codificación diario durante las iteraciones, permitiendo una
planificación granular y una asignación clara de responsabilidades dentro del equipo de desarrollo.

A continuación, se presenta un ejemplo de las \ac{TI} correspondientes a las \ac{HU} definidas, el resto se puede
encontrar en el anexo \ref{anx:ti}.


\begin{engineeringtask}[tb:engtask_t020]
\engtaskuserstory{4}
\engtaskname{Exponer endpoints de ciclo operativo de FM}
\engtasktype{Crear controlador}
\engtaskpointestimation{8}
\engtaskstartdate{14}{5}{2026}
\engtaskenddate{15}{5}{2026}
\engtaskdescription{Implementar el controlador de modelos de características para alta, consulta, actualización y activación/desactivación, con soporte de filtros por dominio, paginación y políticas de autorización según rol funcional.}
\engtaskprogrammer{Jose Luis Echemendia López}
\end{engineeringtask}


\section{Despliegue del sistema}

El despliegue constituye la fase del ciclo de vida del software en la que el sistema 
desarrollado se pone a disposición de los usuarios finales en un entorno de ejecución real o 
de pruebas. Esta etapa implica no solo la instalación y configuración de los componentes 
software en los servidores correspondientes, sino también la verificación de que todas las 
dependencias tecnológicas (bases de datos, sistemas de almacenamiento, colas de mensajes, 
etc.) se encuentren operativas y correctamente parametrizadas \citep{bass2012software}. 
Un despliegue exitoso garantiza que el sistema se comporte en producción de acuerdo con 
las expectativas definidas durante las primeras fases del desarrollo de software, y constituye el paso 
previo indispensable para la ejecución de las pruebas de aceptación por parte del cliente.

\subsection{Diagrama de despliegue}

El diagrama de despliegue es un artefacto de modelado utilizado en el \ac{UML} 
para representar la arquitectura física de un sistema, mostrando 
la disposición de los nodos de hardware, los componentes software que se ejecutan en cada 
uno de ellos y los canales de comunicación que los interconectan \citep{booch2005uml}. Su 
propósito es visualizar la topología de ejecución del sistema, identificar posibles cuellos 
de botella en la comunicación entre nodos y planificar la distribución de los componentes 
en entornos productivos. La figura ilustra el diagrama de despliegue 
correspondiente a la propuesta de solucicón desarrollada.

\section{Fase de pruebas}

La prueba es el proceso de ejecución de un programa con la intención de encontrar errores.
Es crucial para validar el correcto funcionamiento del sistema y asegurar su calidad.
Esta sección presenta la estrategia general de pruebas y los resultados obtenidos al aplicar un conjunto de
técnicas de verificación. El objetivo es demostrar que la aplicación cumple con los requisitos funcionales y
no funcionales definidos, es confiable bajo diferentes condiciones de operación.

\subsection{Estrategia de pruebas}

Una estrategia de pruebas es un plan de alto nivel que define el enfoque, los objetivos, los recursos y el
calendario para las actividades de evaluación de la calidad del software. Su propósito es garantizar que las
pruebas sean sistemáticas, eficientes y efectivas para identificar defectos. \citep{garcia2023}

Se adoptará una estrategia de pruebas basada en riesgos y requisitos. La ejecución seguirá un orden
progresivo: comenzando con las Pruebas Unitarias para cimentar la calidad del código, seguidas de las
Pruebas de Integración para verificar la comunicación entre módulos. En paralelo, se realizarán Pruebas
de Seguridad para validar los controles de acceso. Posteriormente, se ejecutarán Pruebas de Rendimiento
para evaluar la estabilidad bajo carga. El ciclo culminará con las Pruebas de Aceptación para validar el cumplimiento de
las necesidades del usuario final. Como parte integral del proceso, al final de cada iteración se ejecutaron
Pruebas de Regresión para verificar que las nuevas funcionalidades no afectaran el comportamiento existente
del sistema. Esta aproximación por capas asegura una cobertura comprehensiva y una detección temprana
de errores en todas las dimensiones de calidad del sistema.

\subsection{Pruebas unitarias}

Las pruebas unitarias son un tipo de prueba de software donde se verifican unidades individuales o
componentes del código fuente (como una función, un método o una clase) de forma aislada. El objetivo es
validar que cada parte del software funciona correctamente por sí misma \citep{bedoyaAlzate2021}. En este
proyecto, se implementaron pruebas unitarias para cubrir los componentes clave del sistema, asegurando que
cada parte funcione según lo esperado bajo diversas condiciones. Los módulos seleccionados para las pruebas unitarias
son los: servicios del sistema, repositorios, excepciones personalizadas y utilidades de validación. 
Estos componentes son fundamentales para la lógica de negocio, la interacción con la base de datos y 
la gestión de errores, por lo que su correcto funcionamiento es crítico para la estabilidad y confiabilidad 
del sistema.

Para cada componente, se definieron casos de prueba que cubren tanto escenarios típicos como condiciones 
límite o errores esperados. Se utilizaron técnicas de aislamiento mediante \textit{mocks} para simular dependencias
externas como la \ac{BD}, permitiendo así verificar el
comportamiento interno sin efectos colaterales ni dependencias del entorno. A continuación, se presenta un
ejemplo de los casos de pruebas unitarias, el resto se puede encontrar en el anexo \ref{anx:pu}.

\begin{longtable}{|p{3.0cm}|p{10.0cm}|}
    \caption{Caso de Prueba Unitaria No1}\\
    \hline
    \multicolumn{2}{c|}{\colorentry{\textbf{Caso de Prueba Unitaria}}} \\
    \hline
    \endfirsthead
                    
	\multicolumn{2}{c}
    {{\bfseries \tablename\ \thetable{} -- Continuación de la página anterior.}} \\
                    
    \hline
	\multicolumn{2}{c|}{\colorentry{\textbf{Caso de Prueba Unitaria}}} \\
    \hline
    \endhead
                    
    \hline
	\multicolumn{2}{|r|}{{Continúa en la siguiente página}} \\
    \endfoot
                    
    \hline
    \endlastfoot

    % Cuerpo de la tabla
    Descripción: & Verificar que \texttt{FeatureModelVersionManager.}
    \texttt{create\_new\_version(source\_version=None)} 
    crea una versión en estado \texttt{DRAFT} con \texttt{version\_number} consecutivo. \\
	\hline
    Entrada: & 
    \begin{itemize}
        \item \texttt{feature\_model.id} válido.
        \item Mock de \texttt{repository.get\_latest\_version\_number()} retornando 3.
        \item \texttt{source\_version=None}.
    \end{itemize}
    \\
	\hline
    Salida esperada: & 
    \begin{itemize}
        \item Nueva versión con \texttt{version\_number=4}.
        \item \texttt{status=DRAFT}.
        \item Se ejecuta \texttt{session.commit()} y \texttt{session.refresh()}.
    \end{itemize}
    \\
	\hline
    Evaluación de la prueba: & 
    \begin{itemize}
        \item Pasa si se cumplen valores y llamadas a sesión.
        \item Falla si no incrementa versión o no persiste cambios.
    \end{itemize}
    \\
	\hline
\end{longtable}

\subsection{Pruebas de integración}

Las pruebas de integración verifican la interacción y comunicación correcta entre diferentes módulos o
componentes de software una vez que han sido integrados. El objetivo es encontrar fallos en las interfaces y
en la interacción entre unidades integradas \citep{paz2016}. En el contexto de este proyecto que es un servidor
\textit{backend}, las pruebas de integración se realizarán sobre los controladores para verificar que 
las rutas, la lógica de negocio y la interacción con los servicios y repositorios funcionen correctamente. 
Se diseñaron casos de prueba que simulen escenarios reales de uso, incluyendo la validación de datos, la 
gestión de errores y la respuesta del sistema ante diferentes tipos de solicitudes.

Los casos de pruebas de integración se aplicarán sobre las rutas críticas del sistema, las cuales son los controladores
asociados a los servicios de \ac{FM}.Se diseñaron 18 casos de prueba que abordan las integraciones más críticas del sistema. 
A continuación, se presenta un ejemplo de las pruebas de integración, el resto se puede encontrar en el anexo \ref{anx:pi}.

\begin{longtable}{|p{3.0cm}|p{10.0cm}|}
\caption{Caso de Prueba de Integración No1}\\
\hline
\multicolumn{2}{c|}{\colorentry{\textbf{Caso de Prueba de Integración}}} \\
\hline
\endfirsthead
                
\multicolumn{2}{c}
{{\bfseries \tablename\ \thetable{} -- Continuación de la página anterior.}} \\
                
\hline
\multicolumn{2}{c|}{\colorentry{\textbf{Caso de Prueba de Integración}}} \\
\hline
\endhead
                
\hline
\multicolumn{2}{|r|}{{Continúa en la siguiente página}} \\
\endfoot
                
\hline
\endlastfoot
% Cuerpo de la tabla
ID de prueba: HU4-P1: & \ac{HU}: 4 \\
\hline
Nombre: & 
Crear \ac{FM} en dominio válido
\\
\hline
Descripción: & 
Verificar que \texttt{POST} \texttt{/feature-models/} crea un \ac{FM} y persiste metadatos correctamente.
\\
\hline
Condiciones de ejecución: & 
\begin{itemize} 
    \item Sistema levantado con datos de prueba
    \item \ac{BD} limpia o controlada
    \item Usuario autenticado con rol \texttt{MODEL\_DESIGNER} o \texttt{ADMIN}
    \item Existe \texttt{domain\_id} activo
\end{itemize}
\\
\hline
Pasos de ejecución: & 
\begin{enumerate}
    \item Enviar \texttt{POST} \texttt{/api/v1/feature-models/} con \texttt{name}, \texttt{description}, \texttt{domain\_id}
    \item Validar \texttt{HTTP} \texttt{201}
    \item Verificar respuesta con \texttt{id}, \texttt{owner\_id}, \texttt{domain\_id}, \texttt{is\_active=true}
    \item Consultar \texttt{GET} \texttt{/api/v1/feature-models/{id}}
    \item Confirmar persistencia en \ac{BD} (registro creado)
\end{enumerate}
\\
\hline
Resultados esperados: &
\begin{itemize}
    \item Se crea el FM correctamente
    \item El \texttt{owner\_id} corresponde al usuario autenticado
    \item Se registra \texttt{created\_at}/\texttt{updated\_at}
\end{itemize}
\\
\hline
Evaluación de la prueba: &
Aprobada \\
\hline
\end{longtable}

\subsection{Pruebas de seguridad}

Las pruebas de seguridad constituyen un conjunto de técnicas de verificación orientadas a identificar 
vulnerabilidades, debilidades y riesgos en un sistema software que puedan ser explotados por actores maliciosos 
\citep{felderer2016security}. A diferencia de las pruebas funcionales, que verifican que el sistema hace 
lo que debe hacer, las pruebas de seguridad buscan aquello que el sistema no debería permitir: accesos no 
autorizados, fugas de información, ejecución de código malicioso o violaciones de las políticas de control 
de acceso \citep{felderer2015survey}. En el contexto de un servicio \textit{backend} que gestiona \ac{FM} que
tienen información curricular sensible y que pudiera integrarse con sistemas institucionales existentes, 
la seguridad se convierte en un atributo de calidad crítico que debe ser validado rigurosamente.

Los principales objetivos de las pruebas de seguridad incluyen: identificar activos digitales que requieren 
protección, clasificar vulnerabilidades potenciales, evaluar las consecuencias de una posible explotación, y 
proporcionar recomendaciones para la remediación de las fallas detectadas \citep{moradov2023security}. 
En este proyecto, las pruebas de seguridad se centran en tres dimensiones fundamentales: la autenticación y 
gestión de sesiones, la autorización basada en roles estáticos, y la protección de los endpoints de la \ac{API} \ac{REST} 
contra ataques comunes (inyección, manipulación de parámetros, exposición de información sensible). 
Esta sección presenta los resultados de la verificación sistemática de los criterios de seguridad
definidos, utilizando para ello una lista de verificación estructurada que permite evaluar de manera objetiva
el cumplimiento de cada uno de los requisitos establecidos. Dicha lista se evidencia en el anexo \ref{anx:lvs}.


\subsection{Pruebas de rendimiento}

Las pruebas de rendimiento evalúan la velocidad, capacidad de respuesta, estabilidad, confiabilidad y
uso de recursos de un sistema bajo una carga de trabajo determinada. Su objetivo no es encontrar defectos
funcionales, sino identificar y eliminar cuellos de botella en el rendimiento \citep{diaz2008}. 
En el contexto del servicio \textit{backend} para la gestión de variabilidad curricular, 
estas pruebas se centran en validar el cumplimiento de los requisitos no funcionales establecidos, particularmente 
aquellos relacionados con los tiempos de respuesta de la \ac{API}, capacidad de manejar alto tráfico y 
el tiempo de latencia entre los repositorios y la \ac{BD}. Esta sección presenta los resultados de la verificación sistemática
de los criterios de rendimiento definidos, utilizando para ello métricas objetivas que permiten evaluar el
cumplimiento de cada uno de los requisitos establecidos.


\subsection{Pruebas de aceptación}

Las pruebas de aceptación constituyen la fase final del proceso de verificación y validación del software, 
en la cual se determina si el sistema satisface los criterios de aceptación establecidos por el cliente y se encuentra 
apto para su despliegue en producción \citep{pressman2014se}. A diferencia de las pruebas funcionales y no funcionales 
ejecutadas por el equipo de desarrollo, las pruebas de aceptación son típicamente realizadas por el cliente o por usuarios 
finales en un entorno controlado que simula las condiciones reales de operación \citep{crispin2008agile}. Su propósito 
fundamental es generar confianza en que el software resuelve efectivamente los problemas del dominio para el cual fue 
concebido y cumple con las expectativas de sus \textit{stakeholders}.

En el contexto del servicio \textit{backend} para la gestión de variabilidad curricular, las pruebas de aceptación se centran en 
validar los criteros de acepcatión definidos en las \ac{HU}, principalmente en aquellas orientas a modelar y representar
fielmente los \ac{FM} y generar configuraciones válidas que correspondan a itinerarios formativos coherentes.
Estas pruebas se ejecutan sobre escenarios de uso reales, utilizando modelos curriculares de asignaturas y carreras. ç
Se consideran exitosas cuando todas las historias de usuario priorizadas en el plan de entregas han sido 
satisfactoriamente validadas por el cliente. A continuación, se presenta un ejemplo de los casos de pruebas definidos, el 
resto se pueden consultar en el anexo \ref{anx:pa}.

\begin{acceptancetest}[tb:HU6_P2]
\testcasecode{HU6\_P2}
\testcasedescription{Validar que el modelo se puede exportar en formato \ac{UVL} válido.}
\testcaseexeccond{
\begin{itemize}
    \item Existe un modelo con versión PUBLISHED.\\
    \item La versión tiene estructura completa.\\
    \item Usuario está autenticado.
\end{itemize}
}
\testcaseexecstep{
\begin{enumerate} 
    \item Enviar \texttt{GET} a \texttt{/api/v1/models/\{model\_id\}/export?format=uvl} con cabeceras.\\
    \item Validar que la respuesta contenga estada 200.\\
    \item Verificar que el \texttt{Content-Type} es text/plain o application/x-uvl.\\
    \item Validar que el contenido es \ac{UVL} válido (comienza con features o namespace, tiene sintaxis correcta).\\
    \item Guardar contenido en archivo .uvl y validar.
\end{enumerate}
}
\testcaseexpresult{
\begin{itemize} 
    \item Estado \ac{HTTP}: \texttt{200 OK}.\\
    \item \texttt{Content-Type} correcto.\\
    \item \ac{UVL} contiene sintaxis válida.\\
    \item Todas las características, grupos y restricciones están representadas.\\
    \item Puede descargarse como archivo.
    \item Re-importación genera estructura idéntica.
\end{itemize}
}
\testcasename{Exportar modelo en formato \ac{UVL}}
\testcaseuserstory{6}
\end{acceptancetest}


\section{Conclusiones del capítulo}


